\lyseisE{Επανάληψη κεφαλαίου 4}
\solr{4}{1} $5x-3$ και $2\nu+1$.
\solr{4}{2} $(7-3)x+(2+5)y=4x+7y$.
\solr{4}{3} α) $6{,}6$ \quad β) $6$ \quad γ) $12{,}5$ \quad δ) $10$ \quad ε) $4{,}5$.
\solr{4}{4} α) $3x=27$, $x=9$ \quad β) $5x=20$, $x=4$ \quad γ) $3x=33$, $x=11$ \quad δ) $\frac{x}{4}=4$, $x=16$.
\solr{4}{5} α) Ταυτότητα. β) Αδύνατη. γ) Μία λύση, $x=0$. δ) Αδύνατη.
\solr{4}{6} $2{,}5\cdot1{,}2\cdot x=9$, $3x=9$, $x=3$ μ.
\solr{4}{7} $5+0{,}03x=9{,}2$, $0{,}03x=4{,}2$, $x=140$ λεπτά.
\solr{4}{8} $\frac{3}{4}x=18$, $x=18\cdot\frac{4}{3}=24$ κομμάτια.
\solr{4}{9} $3x+3=87$, $x=28$: $28,\,29,\,30$.
\solr{4}{10} $x+(x+12)=100$, $x=44$ (η Αγγελική) και $56$ (ο Νίκος).
\solr{4}{11} $\frac{15\cdot14}{2}=105$ χειραψίες.
\solr{4}{12} Η αμοιβή είναι $5x$ και $\frac{3}{5}\cdot5x=3x=240$, $x=80$ € ανά ημέρα.
\solr{4}{13} α) Σωστό. β) Σωστό. γ) Λάθος: $3x+4x=7x$. δ) Σωστό. ε) Λάθος: $x=12-5=7$. στ) Λάθος: κάθε αριθμός την επαληθεύει.
\solr{4}{14} α) Σωστό. β) Λάθος: επαληθεύω και στο πρόβλημα. γ) Λάθος. δ) Λάθος: $x=4$. ε) Σωστό.
